\scanpage{538}
% E538-01
\mspage{572}

% E538-02
\uline{Proposition} \quad 1°) The image $\delta_\rho(Z_\rho)$ contains
$\gamma'_{\rho,\sigma}$, and is
% E538-03
\deleted{\unsure{eq}}
% E538-04
reduced to $\gamma'_{\rho,\sigma}$ \unsure{iff|iff}
$\rho$ and $\rho^{-1}$ not conjugate in $\mathcal{G}$ ---
therefore it is equal to $\gamma_{\rho,\sigma}$
i.e.
% E538-05
$\delta_\rho:\diagram{\begin{tikzcd}[ampersand replacement=\&,column sep=large]
Z_\rho\arrow[r]\&\gamma_{\rho,\sigma}
\end{tikzcd}}$
% E538-06
epimorphism \unsure{iff|iff} $\rho$ and $\rho^{-1}$ are conjugate in $\mathcal{G}$.

% E538-07
2°) The image $\delta_\sigma(Z_\sigma)$ contains
$\gamma''_{\rho,\sigma}$, and it is reduced to
$\gamma''_{\rho,\sigma}$ \unsure{iff|iff} $\sigma$ and
$\sigma^{-1}=-\sigma$ not conjugate in $\mathcal{G}$ ---
\unsure{therefore|otherwise} it is equal to $\gamma_{\rho,\sigma}$ i.e.
% E538-08
$\delta_\sigma:\diagram{\begin{tikzcd}[ampersand replacement=\&,column sep=large]
Z_\sigma\arrow[r]\&\gamma_{\rho,\sigma}
\end{tikzcd}}$
% E538-09
epimorphism \unsure{iff|iff} $\sigma$ and
$\unsure{-\sigma^{-1}|\sigma^{-1}}=-\sigma$ are conjugate in $\mathcal{G}$.

% E538-10
\uline{Corollary 1}. If $\mathcal{G}$ contains as a quotient
% E538-11
\[
SL(2,\mathbb{Z}/3\mathbb{Z})',\quad\text{with}\quad
\rho=\begin{pmatrix}0&1\\-1&1\end{pmatrix},\qquad
\sigma=\begin{pmatrix}0&-1\\1&0\end{pmatrix}
\]
% E538-12
$(\bmod\ \pm1)$, then the image of $\delta_\rho$ is
$\gamma'_{\rho,\sigma}$.

% E538-13
Indeed, we see that $\rho$ and $\pm\rho^{-1}$ are not
conjugate in $\mathcal{G}$, for they are
not in $SL(2,\mathbb{Z}/3\mathbb{Z})'$ ---
% E538-14
\deleted{\unsure{for example}\ \unsure{for}}
% E538-15
\inserted{Let us check this:}
% E538-16
$\rho$ and $-\rho^{-1}$ already have distinct traces
$1$ and $-1$, $\unsure{\in}(\mathbb{Z}/3\mathbb{Z})^{\unsure{*}}$,
they are not conjugate.

% E538-17
We have
% E538-18
\[
(\underbrace{\tau\sigma}_{\tau'_\infty})(\rho)=\rho^{-1},
\]
% E538-19
therefore
% E538-26
\deleted{all} the elements of
% E538-20
\deleted{$f\in$}
% E538-21
$GL(2,\mathbb{Z}/3\mathbb{Z})$
% E538-22
\deleted{such that}
% E538-23
which conjugate $\rho$ to $\rho^{-1}$ are of the form
$f\tau'_\infty$, where
% E538-24
$f=\unsure{c|b}\rho+d$, we \unsure{want}
$\det(f\tau'_\infty)=1$ i.e. $\det(f)=-1$
$(\text{since }\det(\tau'_\infty)=-1)$,
i.e. $c^2+\unsure{cd}+d^2=-1$, impossible.
% E538-25
\margin{537}

% E538-F01
\margin{\begin{center}
537\\
\texttt{https://grothendieck.umontpellier.fr}\\
Université de Montpellier, UMR CNRS 5149, Case courrier 051, Place Eugène Bataillon\\
34095 Montpellier cedex 5 - France
\end{center}}

\scanpage{539}

% E539-01
\uline{Corollary 2}\quad If $\sigma^2=1$, then
% E539-02
$\diagram{\begin{tikzcd}[ampersand replacement=\&,column sep=large]
Z_\sigma\arrow[r]\&\gamma_{\rho,\sigma}
\end{tikzcd}}$
% E539-03
is surjective. If On the other hand $\mathcal{G}$ contains as a quotient
% E539-04
\deleted{$\mathbb{Z}/4\mathbb{Z}$}
% E539-05
$SL(2,\mathbb{Z}/4\mathbb{Z})$
$(\uline{\text{without}}\ \prime\ !)$, with $\rho,\sigma$ as
\unsure{usual}, then
% E539-06
\[
\operatorname{Im}(\delta_\sigma)=\gamma''_{\rho,\sigma}.
\]

% E539-07
We must see that $\sigma$ and $\sigma^{-1}=-\sigma$ not conjugate
in $SL(2,\mathbb{Z}/4\mathbb{Z})$, using that
$\tau_\infty(\sigma)=\sigma^{-1}$, therefore
% E539-08
\deleted{it f}
% E539-09
the \unsure{elements|elements} of $GL(2,\mathbb{Z}/4\mathbb{Z})$ which
conjugate $\sigma$ to $\sigma^{-1}$ are the
% E539-10
\inserted{$g=f\tau_\infty$ with}
% E539-11
$f=t\sigma+\unsure{u|n}$,
% E539-12
\deleted{\unsure{with}}
% E539-13
and therefore that $\det(g)=1$ means $\det f=-1$
i.e. $t^2+\unsure{u|n}^2=-1$, impossible.

% E539-14
\deleted{
Let us set
% E539-15
\[
(34)\qquad
\left\{
\begin{aligned}
SZ_\rho&=
% E539-16
\inserted{\operatorname{Ker}\delta_\rho=}
Z_\rho\cap\mathcal{G}=\operatorname{Centr}_{\mathcal{G}}(\rho).\\
SZ_\sigma&=
% E539-17
\inserted{\operatorname{Ker}\delta_\sigma=}
Z_\sigma\cap\mathcal{G}=\operatorname{Centr}_{\mathcal{G}}(\sigma)\\
SB_{\rho,\sigma}&=\operatorname{Ker}\delta_B
=B_{\rho,\sigma}\cap\mathcal{G}=L_0.
\end{aligned}
\right.
\]
}

% E539-18
\deleted{
% E539-19
\[
(35)\qquad
\begin{aligned}
&S_0\unsure{\mathcal{M}|M}_{\rho,\sigma}
\subset B_{\rho,\sigma}\times\unsure{Z_\rho|N_\rho}
\times Z_\sigma\times\unsure{G_{\rho,\sigma}|\gamma_{\rho,\sigma}}\\
&\text{''}\\
&\bigl\{(u,a,b,U\unsure{\}|)}\ \bigm|\,
\delta_B(u)=\delta_G(U)=\delta_\sigma(b)
=\varepsilon_3(\mu)\delta_\rho(a)\bigr\}
\end{aligned}
\]
}
% E539-20
\margin{538}

% E539-F01
\margin{\begin{center}
538\\
\texttt{https://grothendieck.umontpellier.fr}\\
Université de Montpellier, UMR CNRS 5149, Case courrier 051, Place Eugène Bataillon\\
34095 Montpellier cedex 5 - France
\end{center}}

\scanpage{540}
% E540-01
\mspage{573}

% E540-02
% E540-03
Let us note that by (2), \inserted{(4)}, we have
% E540-04
\[
(35)\qquad \tau_B \overset{\mathrm{dfn}}{=}(\tau,-1)\in B_{\rho,\sigma}
\]
% E540-05
And obviously $\tau_B\notin L_0=B_{\rho,\sigma}\cap\mathcal{G}$.
So we have an embedding
% E540-06
\[
(36)\qquad
\left\{
\begin{array}{l}
\diagram{\begin{tikzcd}[ampersand replacement=\&,column sep=large,row sep=small]
\mathcal{G}'=\{1,\tau\}.\mathcal{G} \arrow[r,hook] \& G_{\rho,\sigma}\\
\tau \arrow[r,mapsto] \& \tau_B\\
g \arrow[r,mapsto] \& g\quad\text{if }g\in\mathcal{G}
\end{tikzcd}}
\end{array}
\right.,
\]
% E540-07
\margin{\uline{N.B.} $\mathcal{G}'$ is a quotient of
$\mathrm{GL}(2,\mathbf{Z})^\wedge$, and $\tau$ corresponds there to
the image of the element
$\tau_\infty=\begin{pmatrix}-1&0\\0&1\end{pmatrix}$
of $\unsure{\mathrm{GL}(2,\mathbf{Z})}$.}

% E540-08
% E540-09
therefore the \unsure{automorphism} $\tau$ of $\mathcal{G}$ will be considered as
induced by \deleted{\unsure{the ele}} $\operatorname{int}(\tau_B)$,
where $\tau_B$ is itself considered as an
element of $G_{\rho,\sigma}$.

% E540-10
% E540-11
\deleted{\illegible{2}} In connection with the equality
% E540-12
\[
\tau_B(\sigma)=\sigma^{-1}\quad(=-\sigma)
\]
% E540-13
% E540-14
% E540-15
% E540-16
% E540-18
% E540-17
which shows that $\tau_B$ normalizes the subgroup
\inserted{$L_\sigma$} \inserted{generated by $\sigma$}
(isomorphic to a quotient of $\mathbf{Z}/4\mathbf{Z}$),
\deleted{\unsure{group}} \inserted{therefore}
\inserted{normalizes $Z_\sigma=\operatorname{Centr}_{G_{\rho,\sigma}}(L_\sigma)$},
\deleted{\unsure{then let us introduce}} let us introduce
% E540-19
\[
(37)\qquad
N_\sigma=\underset{(\text{$1/2$ direct})}{\{1,\tau_B\}.Z_\sigma},
\]
% E540-20
and the canonical homomorphism
% E540-21
\[
(38)\qquad
\left\{
\begin{array}{l}
\diagram{\begin{tikzcd}[ampersand replacement=\&,column sep=large,row sep=small]
N_\sigma \arrow[r,"i_\sigma"] \& G_{\rho,\sigma} \& \&\\
\tau_B \arrow[r,mapsto] \& \tau_B, \& x \arrow[r,mapsto] \& x\quad\text{if }x\in Z_\sigma
\end{tikzcd}}
\end{array}
\right..
\]

% E540-22
\deleted{This homomorphism is injective iff}
% E540-23
% E540-24
The image of this homomorphism is none other than the normalizer of $\sigma$
in $G_{\rho,\sigma}$, for a \deleted{\unsure{automorphism}}
$g\in G_{\rho,\sigma}$ which normalizes $L_\sigma$ must transform
$\sigma$ into a generator of $L_\sigma$, i.e. into $\sigma$ or
$\sigma^{-1}$, in the first case we will have $g\in Z_\sigma$, in
the second $g(\sigma)=\tau_B(\sigma)$ i.e. $\tau_B^{-1}g\in Z_\sigma$
i.e. $g\in\tau_B Z_\sigma$. The homomorphism (38) is injective

% E540-25
\margin{539}

% E540-F01
\margin{\begin{center}
539\\
\texttt{https://grothendieck.umontpellier.fr}\\
Université de Montpellier, UMR CNRS 5149, Case courrier 051, Place Eugène Bataillon\\
34095 Montpellier cedex 5 - France
\end{center}}

\scanpage{541}
% E541-01
\deleted{iff $\tau_B\in Z$}
% E541-02
iff $\tau_B\notin Z_\sigma$ i.e. $\sigma\ne\sigma^{-1}$ i.e.
$\sigma^2\ne1$ i.e. ``$-1$'' $\ne1$.

% E541-03
While we are at it, let us take advantage of the formula
% E541-04
\[
\tau_B(\rho)=\sigma(\rho^{-1})
\]
% E541-05
to introduce
% E541-06
\[
(39)\qquad
\tau'\uuline{=}\sigma^{-1}\tau\in\mathcal{G}'\subset G_{\rho,\sigma}
\]
% E541-07
(which corresponds to the element
% E541-08
% E541-09
\[
(40)\qquad
\tau'_\infty=\sigma_\infty^{-1}\tau_\infty
=-\sigma_\infty\tau_\infty
=\begin{pmatrix}0&1\\1&0\end{pmatrix}
\in\underbrace{\mathrm{GL}(2,\mathbf{Z})}_{
\inserted{\substack{\Vert\\\mathcal{T}_{1,1}}}}
\biggr),
\]

% E541-10
% E541-11
\deleted{\unsure{which is}} we will therefore have
% E541-12
\[
(40)\qquad \tau'^2=1,\qquad \tau'(\rho)=\rho^{-1}.
\]

% E541-13 deletion
\deleted{Let us therefore introduce}
% E541-14 deletion
\deleted{$(41)\qquad N_\rho=\{1,\tau'\}.$}
% E541-15
% E541-16
\inserted{Thus $\tau'$ normalizes the subgroup $L_\rho$
(isomorphic to a quotient of $\mathbf{Z}/6\mathbf{Z}$) generated
by $\rho$, therefore it normalizes the centralizer
\inserted{$Z_\rho$} of $L_\rho$, and we can construct}
% E541-17
% E541-18
\margin{\uline{N.B.} We have at the same time
$\tau'(\sigma)=\sigma^{-1}$, therefore we can consider $\tau'$
as an element of $N_\sigma$, and \deleted{\illegible{1}} we have
$N_\sigma=\{1,\tau'\}.Z_\sigma$ (which is $1/2$ direct)}
% E541-19
\[
(41)\qquad
N_\rho\overset{\mathrm{dfn}}{=}
\underset{(\text{semidirect})}{\{1,\tau'\}.Z_\rho}
\]
% E541-20
% E541-21
% E541-22
and \deleted{\unsure{a}} \deleted{\unsure{of the}} canonical homomorphism
% E541-23
\[
(42)\qquad
\left\{
\begin{array}{l}
\diagram{\begin{tikzcd}[ampersand replacement=\&,column sep=large,row sep=small]
N_\rho \arrow[r,"i_\rho"] \& G_{\rho,\sigma} \& \&\\
\tau' \arrow[r,mapsto] \& \tau', \& x \arrow[r,mapsto] \& x\quad\text{for }x\in Z_\rho
\end{tikzcd}}
\end{array}
\right..
\]

% E541-24
% E541-25
% E541-26
% E541-27
The image of this homomorphism is the normalizer of $L_\rho$ in
$G_{\rho,\sigma}$ --- we see this as \unsure{a moment ago} for $i_\sigma$,
noting that \deleted{the}
\inserted{$\rho,\rho^{-1}$ are the \uline{only}} generators of
$L_\rho$. The homomorphism $i_\rho$ is injective iff $\tau'\notin Z_\rho$
i.e. \deleted{\unsure{iff}} $\rho\ne\rho^{-1}$ i.e. $\rho^2\ne1$
i.e. $\rho\ne-1$
% E541-28
% E541-29
% E541-30
% E541-31
% E541-32
(\deleted{\unsure{in this case}}
\inserted{N.B. \uline{$\rho=-1$}}
\deleted{\unsure{implies}} \deleted{\unsure{equivalent}}
\deleted{\unsure{in fact}} that $\rho$ commutes with $\sigma$, and that
$\mathcal{G}$ is a quotient of $\mathbf{Z}/12\mathbf{Z}$, with
$\rho$ corresponding to $2$
% E541-33 insertion
\inserted{and $\sigma$ to $3$, but here $2\rho=0$ therefore
$\mathcal{G}$ is a quotient of $\mathbf{Z}/4\mathbf{Z}$ with
$\rho$ corresponding to $2$, $\sigma$ to $3=-1\ldots$}).

% E541-34
% E541-35
\margin{
\[
(43)\qquad\left\{
\begin{array}{l}
\diagram{\begin{tikzcd}[ampersand replacement=\&,column sep=huge,row sep=large]
N_\rho \arrow[r,"\delta_\rho=\delta_G\circ i_\rho"] \& \gamma_{\rho,\sigma}\\
N_\sigma \arrow[r,"\delta_\sigma=\delta_G\circ i_\sigma"] \& \gamma_{\rho,\sigma}
\end{tikzcd}}\\
\hfill\text{epimorphisms}
\end{array}
\right.
\]}
% E541-36
\margin{``We have done what was needed'' to have
\uline{epimorphisms} $\unsure{\varphi}$.}
% E541-37
\margin{540}

% E541-F01
\margin{\begin{center}
540\\
\texttt{https://grothendieck.umontpellier.fr}\\
Université de Montpellier, UMR CNRS 5149, Case courrier 051, Place Eugène Bataillon\\
34095 Montpellier cedex 5 - France
\end{center}}